\label{cold}\section{Cold Plasma Model}
In this section, we use the cold plasma model to calculate the wave
properties. This approximation is valid only if the following conditions are satisfied:

\begin{itemize}
\item  The wave must exist in the cold plasma model, which is not the case of
electrons Bernstein waves (EBW). The description of EBW requires a hot plasma model.

\item  The wave remains far from resonances, where $k_{\bot}\rightarrow\infty
$. Near resonances, mode conversion to electrostatic waves (IBW, EBW) occur,
which are not described by cold plasma wave theory. LHCD and ECCD usually take
place far enough from the LH and UH resonances, respectively, so that the cold
plasma description is valid.

\item  Hot plasma effects, such as the energy flow carried by the coherent
motion of particles $\mathbf{\Phi}_{bT}$, must be negligible.

\item  FLR effects must be small, that is, we must have%
\begin{equation}
\left|  z\right|  =\left|  \frac{k_{\perp}v_{\perp}}{\Omega}\right|  \ll1
\end{equation}
This condition generally holds for LHW and ECW as long as the temperature is
not too high ($T\leq10$ keV) and the resonance harmonic is low ($n=0,1,2$).
\end{itemize}

In the first subsection, we discuss the modeling of RF $k_{\Vert}$ spectrum.
In the second subsection, we use the cold plasma model to calculate wave
properties. In the third and fourth subsections, we use further approximations
to give analytical formulas for the properties of LHW and ECW. This leads to a
simplified description of LHCD and ECCD and allows us to compare our models
and results with other codes.

\subsection{Wave Equation and Dispersion Tensor}

We consider the two-fluids description of a non relativistic plasma in a
constant magnetic field $\mathbf{B}_{0}=B_{0}\widehat{z}$. In a continuous
homogeneous linear medium, a Fourier component $\mathbf{E}_{b}\left(
\mathbf{k}_{b},\omega_{b}\right)  $ verifies the wave equation \cite{sti92b1}%

\begin{equation}
\mathbf{N}_{b}\times(\mathbf{N}_{b}\times\mathbf{E}_{b})+\mathbb{K}%
\cdot\mathbf{E}_{b}=0
\end{equation}
or equivalently
\begin{equation}
(\mathbf{N}_{b}\cdot\mathbf{E}_{b})\mathbf{N}_{b}-N_{b}^{2}\mathbf{E}%
_{b}+\mathbb{K}\cdot\mathbf{E}_{b}=0
\end{equation}
where $\mathbf{N}_{b}=c\mathbf{k}_{b}/\omega_{b}$ is the wave refractive index
and $\mathbb{K}$ is the dielectric tensor. This equation can be expressed in
tensorial form as
\begin{equation}
\mathbb{D}\cdot\mathbf{E}_{b}=0 \label{waveeq3}%
\end{equation}
where the dispersion tensor is
\begin{equation}
\mathbb{D=}\mathbf{N}_{b}\mathbf{N}_{b}\mathbf{-}N_{b}^{2}\mathbb{I}%
+\mathbb{K}%
\end{equation}

By evaluating the susceptibilities in the cold plasma limit, the following
dielectric tensor is obtained \cite{bon84} in the frame $\left(  x,y,z\right)
$:
\begin{equation}
\mathbb{K}=%
\begin{pmatrix}
K_{\perp} & -iK_{xy} & 0\\
iK_{xy} & K_{\perp} & 0\\
0 & 0 & K_{\Vert}%
\end{pmatrix}
\label{ktensor}%
\end{equation}
whose elements in a plasma made of species $s$ are given by
\begin{align}
K_{\perp}  &  =1-\sum_{s}\frac{\omega_{ps}^{2}}{\omega^{2}-\omega_{cs}^{2}%
}\nonumber\\
K_{\parallel}  &  =1-\sum_{s}\frac{\omega_{ps}^{2}}{\omega^{2}}\nonumber\\
K_{xy}  &  =-\sum_{s}\frac{\omega_{ps}^{2}\omega_{cs}}{\omega(\omega
^{2}-\omega_{cs}^{2})}%
\end{align}
where
\begin{equation}
\omega_{ps}=\sqrt{\frac{n_{s}q_{s}^{2}}{\varepsilon_{0}m_{s}}}%
\end{equation}
is the plasma frequency and
\begin{equation}
\omega_{cs}=\frac{\left|  q_{s}\right|  B_{0}}{m_{s}}%
\end{equation}
the cyclotron frequency for the species $s$.

\subsection{Dispersion Relation}

In order to have a non-trivial solution to the wave equation, the dispersion
relation must therefore be satisfied:
\begin{equation}
D\left(  \mathbf{k}_{b},\omega_{b}\right)  =\left|  \mathbb{D}\right|  =0
\label{disprel}%
\end{equation}

The axis $\widehat{x}$ is chosen such that
\begin{equation}
\mathbf{N}_{b}\mathbf{=}N_{b\perp}\widehat{x}+N_{b\parallel}\widehat{z}%
\end{equation}

Taking the parallel component $N_{b\parallel}$ as given, (\ref{disprel}) lead
to the following equation for $N_{b\perp}$%
\begin{equation}
AN_{b\perp}^{4}+BN_{b\perp}^{2}+C=0
\end{equation}
with
\begin{align}
A  &  =K_{\perp}\\
B  &  =\left(  N_{b\parallel}^{2}-K_{\perp}\right)  \left(  K_{\perp
}+K_{\parallel}\right)  +K_{xy}^{2}\\
C  &  =K_{\parallel}\left[  \left(  N_{b\parallel}^{2}-K_{\perp}\right)
^{2}-K_{xy}^{2}\right]
\end{align}
or equivalently
\begin{align}
A  &  =K_{\perp}\\
B  &  =\left(  K_{\perp}+K_{\parallel}\right)  N_{b\parallel}^{2}-\left(
K_{R}K_{L}+K_{\perp}K_{\parallel}\right) \\
C  &  =K_{b\parallel}\left(  N_{b\parallel}^{2}-K_{R}\right)  \left(
N_{b\parallel}^{2}-K_{L}\right)
\end{align}
with
\begin{align}
K_{R}  &  =K_{\perp}+K_{xy}\\
K_{L}  &  =K_{\perp}-K_{xy}%
\end{align}

We find the expression for $N_{b\perp}$ as a function of $N_{b\parallel}$ and
$\omega$:
\begin{equation}
N_{b\perp}^{2}=\frac{-B\pm\sqrt{B^{2}-4AC}}{2A} \label{nperpform}%
\end{equation}

We see immediately that the resonances occur for $A=0$, that is,
\begin{equation}
K_{\perp}=0
\end{equation}
and that cut-offs occur for $C=0$, that is,
\begin{equation}
K_{R}=N_{b\parallel}^{2}%
\end{equation}%
\begin{equation}
K_{L}=N_{b\parallel}^{2}%
\end{equation}
or
\begin{equation}
K_{\parallel}=0
\end{equation}

\subsection{Polarization components}

Once $N_{b\perp}$ has been evaluated, the components of the polarization of
the electric field are simply given as the eigenvector of the wave equation
(\ref{waveeq3})
\begin{equation}
\mathbb{D}\cdot\mathbf{E}_{b}=0 \label{polform}%
\end{equation}

\subsection{Power flow}

The relation between power flux and electric field was described using the
vector $\mathbf{\Phi}_{b}$ defined in (\ref{phib2})
\begin{equation}
\mathbf{\Phi}_{b}=\mathbf{\Phi}_{bP}+\mathbf{\Phi}_{bT}%
\end{equation}
with (\ref{phibp2})
\begin{equation}
\mathbf{\Phi}_{bP}=\text{Re}\left[  \mathbf{N}_{b}-\left(  \mathbf{N}_{b}%
\cdot\mathbf{e}_{b}\right)  \mathbf{e}_{b}^{\ast}\right]  \label{phibp}%
\end{equation}
and (\ref{phit})
\begin{equation}
\mathbf{\Phi}_{bT}=\mathbf{-}\frac{1}{2}\mathbf{e}_{b}^{\ast}\cdot
\frac{\partial\mathbb{K}^{H}}{\partial\mathbf{N}_{b}}\cdot\mathbf{e}_{b}%
\end{equation}

In the cold plasma model, the dielectric tensor $\mathbb{K}$ is independent of
$\mathbf{N}_{b}$ and the contribution $\mathbf{\Phi}_{bT}$, from the coherent
motion of particles, vanishes.%
\begin{equation}
\mathbf{\Phi}_{bT}=0
\end{equation}
It is usually a very good approximation to neglect this kinetic power flux,
even for quasi-electrostatic waves like Lower-Hybrid waves (See \ref{kinetic}).

\subsection{Conclusion}

The cold plasma model gives expressions (\ref{nperpform}), (\ref{polform}) and
(\ref{phibp}) for the calculation of $N_{\bot}$, $e_{b}$ and $\mathbf{\Phi
}_{b}$ respectively. These formulas can be generally used for the calculation
of the diffusion coefficient in LHCD and ECCD.

\section{Lower Hybrid Current Drive\label{LHCD}}

Analytic expression for the properties of LH waves can be obtained in the cold
plasma description if further approximations are made. Then, an analytical
expression can be obtained for the LH diffusion coefficient.

\subsection{Electrostatic Dispersion Relation}

We consider the wave equation (\ref{waveeq3})
\begin{equation}
(\mathbf{N}_{b}\cdot\mathbf{E}_{b})\mathbf{N}_{b}-N_{b}^{2}\mathbf{E}%
_{b}+\mathbb{K}\cdot\mathbf{E}_{b}=0 \label{waveeq}%
\end{equation}

The electric field can be separated into its longitudinal and transverse
components with respect to the normalized wave vector $\mathbf{N}_{b}$:
\begin{equation}
\mathbf{E}_{b}=\mathbf{E}_{bL}+\mathbf{E}_{bT} \label{edecomp}%
\end{equation}
so that the wave equation (\ref{waveeq}) becomes
\begin{equation}
(N_{b}^{2}-\mathbb{K})\mathbf{E}_{bT}-\mathbb{K}\cdot\mathbf{E}_{bL}=0
\label{waveeq2}%
\end{equation}

Electrostatic waves, and also electromagnetic waves approaching resonances,
such as LHW, can satisfy the condition
\begin{equation}
N_{b}^{2}\gg|K_{ij}| \label{esapprox}%
\end{equation}
and therefore (\ref{waveeq2}) reduces to
\begin{equation}
N_{b}^{2}\mathbf{E}_{bT}-\mathbb{K}\cdot\mathbf{E}_{bL}\simeq0
\label{waveeqes}%
\end{equation}

Dot-multiplying (\ref{waveeqes}) by $\widehat{N}_{b}=\mathbf{N}_{b}/N_{b}$
leads to the electrostatic dispersion relation
\begin{equation}
D_{L}\equiv\widehat{N}_{b}\cdot\mathbb{K}\cdot\widehat{N}_{b}=0
\label{dispreles}%
\end{equation}
and the transverse electric field is given by
\begin{equation}
\mathbf{E}_{bT}=\frac{1}{N_{b}^{2}}\mathbb{K}\cdot\mathbf{E}_{bL} \label{et}%
\end{equation}

We note from (\ref{et}) and (\ref{esapprox}) that $|\mathbf{E}_{bT}%
|\ll|\mathbf{E}_{bL}|$. The electric field is quasi longitudinal, which
justify the term electrostatic approximation given to the condition
(\ref{esapprox}).

\subsection{Cold Plasma Limit}

Using (\ref{ktensor}), the electrostatic dispersion relation (\ref{dispreles})
in the cold plasma limit is then given by
\begin{equation}
D_{L}=N_{b\perp}^{2}K_{\perp}+N_{b\parallel}^{2}K_{\parallel}=0
\label{disprelescpl}%
\end{equation}
which gives an expression for $N_{b\perp}$ as a function of $N_{b\parallel}$
and $\omega_{b}$:
\begin{equation}
N_{b\perp}^{2}=\frac{-K_{\parallel}}{K_{\perp}}N_{b\parallel}^{2}
\label{nperplh0}%
\end{equation}

\subsection{Lower Hybrid Waves}

We consider the lower-hybrid range of frequency in a electron-ion plasma,
where the following ordering applies
\begin{equation}
\omega_{ci}\ll\omega_{b}\ll\omega_{ce}%
\end{equation}

Assuming in addition that
\begin{equation}
\omega_{b}\ll\omega_{pe}%
\end{equation}
leads to approximate expressions of the dielectric tensor components given
by:
\begin{align}
K_{\perp}  &  \simeq1+\frac{\omega_{pe}^{2}}{\omega_{ce}^{2}}-\frac
{\omega_{pi}^{2}}{\omega_{b}^{2}}\nonumber\\
K_{\parallel}  &  \simeq-\frac{\omega_{pe}^{2}}{\omega_{b}^{2}}\label{kpar}\\
K_{xy}  &  \simeq\frac{\omega_{pe}^{2}}{\omega_{b}\omega_{ce}}
\label{disptenlh}%
\end{align}
where $K_{\perp}$ can be rewritten as%
\begin{align}
K_{\perp}  &  =\left(  1+\frac{\omega_{pe}^{2}}{\omega_{ce}^{2}}\right)
\left(  1-\frac{\omega_{pi}^{2}/\omega_{b}^{2}}{1+\omega_{pe}^{2}/\omega
_{ce}^{2}}\right) \label{kperp}\\
&  =\left(  1+\frac{\omega_{pe}^{2}}{\omega_{ce}^{2}}\right)  \left(
1-\frac{\omega_{\text{LH}}^{2}}{\omega_{b}^{2}}\right) \\
&  =\frac{\omega_{pi}^{2}}{\omega_{b}^{2}}\left(  \frac{\omega_{b}^{2}}%
{\omega_{\text{LH}}^{2}}-1\right)
\end{align}
where
\begin{equation}
\omega_{\text{LH}}^{2}=\frac{\omega_{pi}^{2}}{1+\omega_{pe}^{2}/\omega
_{ce}^{2}}%
\end{equation}
is the lower hybrid frequency.

The perpendicular index of refraction becomes
\begin{equation}
N_{b\perp}^{2}=\frac{\omega_{pe}^{2}/\omega_{pi}^{2}}{\left(  \omega_{b}%
^{2}/\omega_{\text{LH}}^{2}-1\right)  }N_{b\parallel}^{2} \label{nperplh2}%
\end{equation}
We recognize that $N_{b\perp}\rightarrow\infty$ as the the lower-hybrid
frequency approaches the wave frequency. However, in most LHCD scenarios, wave
frequency are sensibly higher than the lower-hybrid frequency in order to
avoid conversion to IBW. Typically, in Alcator C-Mod, we have
\begin{equation}
\frac{\omega_{b}}{\omega_{\text{LH}}}\sim2\text{ or }3
\end{equation}

In that case, and as shown in \ref{kinetic}, the contribution of
$\mathbf{\Phi}_{bT}$ to the power flow can be neglected. The cold plasma
description therefore remains valid.

\subsection{Polarization}

The LH Waves of concern for LHCD are quasi-electrostatic, and therefore the
electric field is quasi-longitudinal ($\mathbf{E}_{b}\parallel\mathbf{N}_{b}$)
and we have
\begin{equation}
e_{b,i}=\frac{E_{b,i}}{\left|  \mathbf{E}_{b}\right|  }\simeq\frac{k_{b,i}%
}{\left|  \mathbf{k}_{b}\right|  }%
\end{equation}
for any component $i$, so that the polarization elements become
\begin{align}
e_{b,+}  &  \equiv\frac{E_{b,x}+iE_{b,y}}{\sqrt{2}\left|  \mathbf{E}%
_{b}\right|  }=\frac{k_{b\bot}}{\sqrt{2}k_{b}}\left(  \cos\alpha+i\sin
\alpha\right)  =\frac{k_{b\bot}}{\sqrt{2}k_{b}}e^{+i\alpha}\label{pollh}\\
e_{b,-}  &  \equiv\frac{E_{b,x}-iE_{b,y}}{\sqrt{2}\left|  \mathbf{E}%
_{b}\right|  }=\frac{k_{b\bot}}{\sqrt{2}k_{b}}\left(  \cos\alpha-i\sin
\alpha\right)  =\frac{k_{b\bot}}{\sqrt{2}k_{b}}e^{-i\alpha}\nonumber\\
e_{b,\Vert}  &  \equiv\frac{E_{b,z}}{\left|  \mathbf{E}_{b}\right|  }%
=\frac{k_{b\parallel}}{k_{b}}\nonumber
\end{align}

\subsection{Determination of $\Theta_{\mathbf{k}}^{b,\text{LH}}$}

Lower Hybrid Current Drive (LHCD) results from momentum exchange from the LH
wave to the plasma through Landau damping (harmonic $n=0$). In this case, the
coefficient (\ref{thetakn2}) becomes
\begin{align}
\Theta_{\mathbf{k}}^{b,0}  &  =\frac{1}{\sqrt{2}}e_{b,+}e^{-i\alpha}%
J_{-1}\left(  \frac{k_{b\bot}v_{\bot}}{\Omega}\right)  +\frac{1}{\sqrt{2}%
}e_{b,-}e^{+i\alpha}J_{1}\left(  \frac{k_{b\bot}v_{\bot}}{\Omega}\right)
+\frac{p_{\Vert}}{p_{\bot}}e_{b,\Vert}J_{0}\left(  \frac{k_{b\bot}v_{\bot}%
}{\Omega}\right) \\
&  =\frac{1}{\sqrt{2}}\left(  e_{b,-}e^{+i\alpha}-e_{b,+}e^{-i\alpha}\right)
J_{1}\left(  \frac{k_{b\bot}v_{\bot}}{\Omega}\right)  +\frac{p_{\Vert}%
}{p_{\bot}}e_{b,\Vert}J_{0}\left(  \frac{k_{b\bot}v_{\bot}}{\Omega}\right)
\end{align}

Using (\ref{pollh}), we see that the perpendicular components cancel and we
are left with
\begin{equation}
\Theta_{\mathbf{k}}^{b,\text{LH}}=\frac{p_{\Vert}}{p_{\bot}}\frac
{k_{b\parallel}}{k_{b}}J_{0}\left(  \frac{k_{b\bot}v_{\bot}}{\Omega}\right)
\label{thetaklh}%
\end{equation}

The argument of the Bessel function in (\ref{thetaklh}) is
\begin{equation}
\frac{k_{b\bot}v_{\bot}}{\Omega}=N_{b\bot}\frac{v_{\bot}}{c}\frac{\omega_{b}%
}{\Omega}%
\end{equation}
where an expression for $N_{\bot}$ is given by (\ref{nperplh2}). The argument
of the Bessel function becomes
\begin{equation}
\frac{k_{b\bot}v_{\bot}}{\Omega}=\frac{1}{\sqrt{\omega_{b}^{2}/\omega
_{\text{LH}}^{2}-1}}\frac{\omega_{pe}}{\omega_{pi}}\frac{\omega_{b}}%
{\omega_{ce}}N_{b\parallel}\frac{p_{\bot}}{p_{Te}}\beta_{Te}%
\end{equation}
and we see that for $\omega_{pe}\sim\omega_{ce}$ and $\omega_{b}\sim
2\omega_{\text{LH}}$, we get
\begin{equation}
\left|  \frac{k_{b\bot}v_{\bot}}{\Omega}\right|  \sim\beta_{Te}N_{b\parallel
}\frac{p_{\bot}}{p_{Te}}%
\end{equation}

Given that most electrons concerned with LHCD have $p_{\bot}\sim p_{Te}$, and
that for LHCD we have typically $N_{b\parallel}\sim2$, we see that it is
reasonable to take the limit
\begin{equation}
\left|  \frac{k_{b\bot}v_{\bot}}{\Omega}\right|  \ll1
\end{equation}
as long as the plasma is not too relativistic ($\beta_{Te}\ll1$). This limit
is consistent with the validity condition of the cold plasma description. In
this limit, valid for most LHCD scenarios, we have
\begin{equation}
\Theta_{\mathbf{k}}^{b,\text{LH}}\simeq\frac{p_{\Vert}}{p_{\bot}}%
\frac{k_{b\parallel}}{k_{b}} \label{thetaklh2}%
\end{equation}

\subsection{Determination of $\mathbf{\Phi}_{bP}^{\text{LH}}$}

The vector $\mathbf{\Phi}_{b}^{\text{LH}}$ describes the relation between the
energy flux and the electric field. Although LHW are almost electrostatic, the
dominant contribution to the power flux is the wave Poynting flux
$\mathbf{\Phi}_{bP}^{\text{LH}}$, assuming that the wave frequency remains
sensibly higher than the LH frequency. The contribution of the Kinetic power
flux is calculated in \ref{kinetic}. In most LHCD scenarios, this contribution
is not more than a few percents of the total flux and can therefore be
neglected. We have (\ref{phibp}).%
\begin{equation}
\mathbf{\Phi}_{bP}=\text{Re}\left[  \left|  e_{bT}\right|  ^{2}\mathbf{N}%
_{b}-N_{b}e_{bL}\mathbf{e}_{bT}^{\ast}\right]  \label{phiblh}%
\end{equation}

LH waves are quasi-electrostatic, so that the longitudinal and transverse
components are given by (\ref{et})
\begin{align}
e_{bL}  &  \simeq1\\
\mathbf{e}_{bT}  &  =\frac{1}{N_{b}^{2}}\mathbb{K}\cdot\mathbf{e}_{bL}%
\end{align}
so that the first term in (\ref{phiblh}) can be neglected, and
\begin{align}
\mathbf{\Phi}_{bP}^{\text{LH}}  &  \simeq\text{Re}\left(  -e_{bL}\frac
{1}{N_{b}}\mathbb{K}^{\ast}\cdot\mathbf{e}_{bL}^{\ast}\right) \nonumber\\
&  =\text{Re}\left(  -\frac{1}{N_{b}}\mathbb{K}^{\ast}\cdot\widehat{N}\right)
\label{poyntinges}%
\end{align}

Note from (\ref{poyntinges}) and (\ref{dispreles}) that the Poynting flux is
in the direction perpendicular to the wave vector:
\begin{equation}
\mathbf{\Phi}_{bP}^{\text{LH}}\cdot\widehat{N}=0
\end{equation}

We finally find, using (\ref{ktensor}),
\begin{equation}
\mathbf{\Phi}_{bP}^{\text{LH}}=-\frac{1}{N_{b}^{2}}\left(  K_{\perp}%
\mathbf{N}_{b\perp}+K_{\Vert}N_{b\Vert}\widehat{z}\right)  \label{phiblh2}%
\end{equation}

\subsection{Determination of $\mathbf{\Phi}_{bT}^{\text{LH}}$ \label{kinetic}}

\bigskip In the analysis above the kinetic part of the power flux, due to the
coherent motion of charge carriers, has been neglected. This approximation
must be justified by comparing the Poynting and the kinetic fluxes.

The normalized expression for the kinetic power flux associated with an
electromagnetic wave in a kinetic plasma is given by (\ref{phit})
\begin{equation}
\mathbf{\Phi}_{bT}=\mathbf{-}\frac{\omega_{b}}{2c}\mathbf{e}_{b}^{\ast}%
\cdot\frac{\partial\mathbb{K}^{H}}{\partial\mathbf{k}}\cdot\mathbf{e}_{b}
\label{kinflux}%
\end{equation}

\subsubsection{Electrostatic Waves}

In the electrostatic approximation, the electric field is quasi-longitudinal
(\ref{et})
\begin{equation}
\left|  \mathbf{e}_{bT}\right|  \ll e_{bL}\simeq1
\end{equation}
and (\ref{kinflux}) becomes
\begin{equation}
\mathbf{\Phi}_{bT}=\mathbf{-}\frac{\omega_{b}}{2c}\widehat{k}_{b}\cdot
\frac{\partial\mathbb{K}^{H}}{\partial\mathbf{k}}\cdot\widehat{k}_{b}
\label{kinfluxes1}%
\end{equation}

In the frame of the wave vector $\widehat{k}_{b}$, (\ref{kinfluxes1}) can then
be expressed as
\begin{equation}
\mathbf{\Phi}_{bT}=\mathbf{-}\frac{\omega_{b}}{2c}\frac{\partial K_{L}^{H}%
}{\partial k}\widehat{k}_{b} \label{kinfluxes2}%
\end{equation}
where
\begin{equation}
K_{L}=\widehat{k}_{b}\cdot\mathbb{K}\cdot\widehat{k}_{b}%
\end{equation}
is the longitudinal (electrostatic) component of the dielectric tensor.

\subsubsection{Lower Hybrid Waves}

An expression of the electrostatic dispersion relation for LH
waves that includes the first-order thermal corrections is given
by \cite{por02n1}:
\begin{equation}
K_{L}^{\text{LH}}(\omega,k_{b})\simeq1+\frac{k_{b\perp}^{2}}{k_{b}^{2}}%
\frac{\omega_{pe}^{2}}{\omega_{ce}^{2}}\left(  1-\frac{3b_{e}}{4}\right)
-\frac{k_{b\perp}^{2}}{k_{b}^{2}}\frac{\omega_{pi}^{2}}{\omega_{b}^{2}}\left(
1+\frac{3b_{i}\omega_{ci}^{2}}{\omega_{b}^{2}}\right)  -\frac{k_{b\parallel
}^{2}}{k_{b}^{2}}\frac{\omega_{pe}^{2}}{\omega_{b}^{2}} \label{dispreleslhkin}%
\end{equation}
where the finite Larmor radius (FRL) effects are scaled by the parameters
\begin{equation}
b_{e}=\frac{k_{b\perp}^{2}v_{Te}^{2}}{\omega_{ce}^{2}},\qquad b_{i}%
=\frac{k_{b\perp}^{2}v_{Ti}^{2}}{\omega_{ci}^{2}} \label{flr}%
\end{equation}
with $v_{T}^{2}=T/m$.

The $k$-dependence of the dispersion relation is found in the thermal
correction terms so that
\begin{equation}
\frac{\partial K_{L}^{\text{LH}}}{\partial k}=-\frac{3}{4}\frac{k_{b\perp}%
^{4}}{k_{b}^{4}}\frac{\omega_{pe}^{2}}{\omega_{ce}^{2}}\frac{2k_{b}v_{Te}^{2}%
}{\omega_{ce}^{2}}-\frac{k_{b\perp}^{4}}{k_{b}^{4}}\frac{\omega_{pi}^{2}%
}{\omega_{b}^{2}}\frac{3\omega_{ci}^{2}}{\omega_{b}^{2}}\frac{2k_{b}v_{Ti}%
^{2}}{\omega_{ci}^{2}} \label{kdepdispreleslhkin}%
\end{equation}

Inside the plasma, we have $k_{b\parallel}^{2}\ll k_{b\perp}^{2}$ and
therefore $k_{b\perp}^{2}/k_{b}^{2}\simeq1$. We finally get the following
expression for the kinetic flux associated with quasi-electrostatic LH waves:
\begin{equation}
\mathbf{\Phi}_{bT}^{\text{LH}}=\frac{1}{N_{b}}\left(  \frac{1}{4}\frac
{\omega_{pe}^{2}k_{b}^{2}v_{Te}^{2}}{\omega_{ce}^{4}}+\frac{\omega_{pi}%
^{2}k_{b}^{2}v_{Ti}^{2}}{\omega_{b}^{4}}\right)  \widehat{k}_{b}
\label{kinfluxeslh}%
\end{equation}

The kinetic flux is oriented in the direction of the wave vector, and the
incident flux on the flux surface
\begin{equation}
\left|  \mathbf{\Phi}_{bT}^{\text{LH}}\cdot\widehat{\psi}\right|  =\frac
{1}{N_{b}}\left(  \frac{1}{4}\frac{\omega_{pe}^{2}k_{b}^{2}v_{Te}^{2}}%
{\omega_{ce}^{4}}+\frac{\omega_{pi}^{2}k_{b}^{2}v_{Ti}^{2}}{\omega_{b}^{4}%
}\right)  \left|  \widehat{k}_{b}\cdot\widehat{\psi}\right|
\label{kinfluxeslhperp}%
\end{equation}
This incident kinetic flux must be compared to the incident Poynting flux
taken in the cold plasma limit from (\ref{phiblh2})
\begin{equation}
\left|  \mathbf{\Phi}_{bP}^{\text{LH}}\cdot\widehat{\psi}\right|
=\frac{K_{\perp}}{N_{b}^{2}}\left|  \mathbf{N}_{b\perp}\cdot\widehat{\psi
}\right|  \simeq\frac{K_{\perp}}{N_{b}}\left|  \widehat{k}_{b}\cdot
\widehat{\psi}\right|  \label{poyntfluxeslhperp}%
\end{equation}

The ratio of the incident kinetic power flux to the Poynting flux is given by
\begin{equation}
\frac{\left|  \mathbf{\Phi}_{bT}^{\text{LH}}\cdot\widehat{\psi}\right|
}{\left|  \mathbf{\Phi}_{bP}^{\text{LH}}\cdot\widehat{\psi}\right|  }=\frac
{1}{K_{\perp}}\left(  \frac{3}{4}\frac{\omega_{pe}^{2}v_{Te}^{2}k_{b}^{2}%
}{\omega_{ce}^{4}}+3\frac{\omega_{pi}^{2}v_{Ti}^{2}k_{b}^{2}}{\omega_{b}^{4}%
}\right)  \label{fluxratio}%
\end{equation}

Using ({\ref{nperplh0}) we have
\begin{equation}
k_{b}^{2}\simeq\frac{\omega_{b}^{2}}{c^{2}}N_{b\perp}^{2}\simeq\frac
{\omega_{b}^{2}}{c^{2}}\frac{\left(  -K_{\parallel}\right)  }{K_{\perp}%
}N_{b\parallel}^{2}\simeq\frac{1}{K_{\perp}}\frac{\omega_{pe}^{2}}{c^{2}%
}N_{b\parallel}^{2}%
\end{equation}
so that
\begin{equation}
\frac{\mathbf{\Phi}_{bT\perp}^{\text{LH}}}{\mathbf{\Phi}_{bP\perp}^{\text{LH}%
}}=\frac{3\beta_{Te}^{2}N_{b\parallel}^{2}}{(1+\omega_{pe}^{2}/\omega_{ce}%
^{2})(1-\omega_{\text{LH}}^{2}/\omega_{b}^{2})}\left(  \frac{1}{4}\frac
{\omega_{pe}^{4}}{\omega_{ce}^{4}}+\frac{\omega_{pi}^{4}}{\omega_{b}^{4}}%
\frac{T_{i}}{T_{e}}\right)  \label{fluxratio2}%
\end{equation}
where $\beta_{Te}^{2}=T_{e}/mc^{2}$. }

It can be seen from (\ref{fluxratio2}) that the kinetic part of the power flux
becomes significant only as the wave approaches the lower-hybrid resonance
very closely. In a typical LHCD context (for instance Alcator C-Mod,
\cite{bon00}), we have
\begin{equation}
\omega_{b}\sim2\omega_{\text{LH}},\ \omega_{pe}\sim\Omega_{e},\ \omega
_{pi}\sim\omega,\ T_{e}\sim T_{i}\text{ and }\beta_{Te}\sim0.1 \label{cmod}%
\end{equation}
so that the kinetic part of the power flux is not more than a few percents.

\subsection{LH Diffusion Coefficient}

\subsubsection{General expression in small FLR\ limit and ES approximation}

The normalized bounce-averaged diffusion coefficient for the Fokker-Planck
equation is given by (\ref{dbnb})%
\begin{align}
\overline{D}_{b}^{\text{LH}(0)}(p,\xi_{0})  &  =\frac{\gamma p_{Te}}{p\left|
\xi_{0}\right|  }\frac{1}{\lambda\widetilde{q}}\frac{r_{\theta_{b}}}{R_{p}%
}\frac{B^{\theta_{b}}}{B_{P}^{\theta_{b}}}\frac{\xi_{0}^{2}}{\xi_{\theta_{b}%
}^{2}}\Psi_{\theta_{b}}\overline{D}_{b,0}^{\text{LH,}\theta_{b}}%
\times\nonumber\\
&  H\left(  \theta_{b}-\theta_{\min}\right)  H\left(  \theta_{\max}-\theta
_{b}\right)  \left[  \frac{1}{2}\sum\limits_{\sigma}\right]  _{T}\delta\left(
N_{b\Vert}-N_{\Vert\text{res}}^{\theta_{b}\text{LH}}\right)  \left|
\Theta_{\mathbf{k,}\theta_{b}}^{b,\text{LH}}\right|  ^{2} \label{dlh2}%
\end{align}
with%
\begin{align}
\overline{D}_{b,0}^{\text{LH,}\theta_{b}}  &  =\frac{1}{r_{\theta_{b}%
}R_{\theta_{b}}}\frac{1}{m_{e}\ln\Lambda}\frac{1}{\omega_{b}\omega_{pe}^{2}%
}\frac{f_{\text{inc},b}^{l+1/2}}{\left|  \mathbf{\Phi}_{b}^{\text{LH}}\right|
}P_{b,\text{inc}}\\
N_{\Vert\text{res}}^{\theta_{b}\text{LH}}  &  =\frac{1}{\beta_{Te}}%
\frac{\gamma p_{Te}}{p\xi_{\theta_{b}}}%
\end{align}

Within the electrostatic approximation and in the small FLR limit, we have
obtained the following expressions for the LHW properties (\ref{thetaklh2}),
(\ref{nperplh0}), (\ref{phiblh2})
\begin{align}
\Theta_{\mathbf{k},\theta_{b}}^{b,\text{LH}}  &  =\frac{\xi_{\theta_{b}}%
}{\sqrt{\Psi_{\theta_{b}}\left(  1-\xi_{0}^{2}\right)  }}\frac{N_{b\parallel}%
}{N_{b}}\\
N_{b\perp}  &  =\sqrt{\frac{-K_{\parallel}}{K_{\perp}}}N_{b\parallel}\\
\mathbf{\Phi}_{b}^{\text{LH}}  &  =-\frac{1}{N_{b}^{2}}\left(  K_{\perp
}\mathbf{N}_{b\perp}+K_{\Vert}N_{b\Vert}\widehat{z}\right)
\end{align}
so that%
\begin{align}
\left|  \mathbf{\Phi}_{b}^{\text{LH}}\right|   &  =\frac{1}{N_{b}^{2}}%
\sqrt{K_{\perp}^{2}N_{b\perp}^{2}+K_{\Vert}^{2}N_{b\Vert}^{2}}\\
&  \simeq\frac{N_{b\Vert}}{N_{b}^{2}}\left|  K_{\Vert}\right| \\
&  =\frac{K_{\perp}}{N_{b\Vert}}%
\end{align}
and where $K_{\perp}$ and $K_{\Vert}$ are given by (\ref{kperp}) and
(\ref{kpar}) respectively

\subsubsection{Simplified expression for LHCD}\label{simplifiedlhcd}

We consider the limit of a large aspect ratio tokamak with circular
flux-surfaces (limit of a cylindrical plasma). In that case, $r_{\theta_{b}%
}/R_{p}\rightarrow0$ and the effects of magnetic trapping disppear. We can use
the following asymptotic expressions
\begin{align}
\Psi_{\theta_{b}}  &  \rightarrow1\\
\lambda\left(  \psi,\xi_{0}\right)   &  \rightarrow1\\
\frac{1}{\widetilde{q}}\frac{r_{\theta_{b}}}{R_{p}}\frac{B^{\theta_{b}}}%
{B_{P}^{\theta_{b}}}  &  \rightarrow1\\
\frac{\xi_{\theta_{b}}}{\xi_{0}}  &  \rightarrow1
\end{align}

Because of cylindrical symmetry, the dependence $\theta_{b}$ on disappears.

The QL diffusion coefficient for LHCD\ (\ref{dlh2}) can therefore be written
as%
\begin{equation}
\overline{D}_{b}^{\text{LH}(0)}(p,\xi_{0})=\frac{\gamma p_{Te}}{p\left|
\xi_{0}\right|  }\frac{\xi_{0}^{2}}{\left(  1-\xi_{0}^{2}\right)  }%
\frac{K_{\perp}}{\left(  -K_{\parallel}\right)  }\overline{D}_{b,0}%
^{\text{LH}}\frac{1}{\Delta N_{b\Vert}}H\left(  N_{\Vert\text{res}}%
^{\text{LH}}-N_{b\Vert\min}\right)  H\left(  N_{b\Vert\max}-N_{\Vert
\text{res}}^{\text{LH}}\right)
\end{equation}
with
\begin{align}
\overline{D}_{b,0}^{\text{LH}}  &  =\frac{1}{rR_{p}}\frac{1}{m_{e}\ln\Lambda
}\frac{1}{\omega_{b}\omega_{pe}^{2}}\frac{N_{b\parallel,0}}{K_{\perp}%
}f_{\text{inc},b}^{l+1/2}P_{b,\text{inc}}\\
N_{\Vert\text{res}}^{\text{LH}}  &  =\frac{1}{\beta_{Te}}\frac{\gamma p_{Te}%
}{p\xi_{0}}%
\end{align}
and where we used $N_{b}\simeq N_{b\perp}$ and assume a square power spectrum
as in (\ref{square}).

In order to compare with LHCD\ operators found in the litterature, we redefine
the LH constant factor such that%
\begin{equation}
\overline{D}_{b}^{\text{LH}(0)}(p,\xi_{0})=\frac{\gamma p_{Te}}{p\left|
\xi_{0}\right|  }\frac{\xi_{0}^{2}}{\left(  1-\xi_{0}^{2}\right)  }%
\overline{D}_{b,0\text{,new}}^{\text{LH}}H\left(  N_{\Vert\text{res}%
}^{\text{LH}}-N_{b\Vert\min}\right)  H\left(  N_{b\Vert\max}-N_{\Vert
\text{res}}^{\text{LH}}\right)
\end{equation}
with
\begin{align}
\overline{D}_{b,0\text{,new}}^{\text{LH}}  &  =\frac{N_{b\parallel,0}^{2}%
}{N_{b\bot,0}^{2}}\frac{1}{\Delta N_{b\Vert}}\overline{D}_{b,0}^{\text{LH}}\\
&  =\frac{1}{rR_{p}}\frac{1}{m_{e}\ln\Lambda}\frac{\omega_{b}}{\omega_{pe}%
^{4}}\frac{N_{b\parallel,0}}{\Delta N_{b\Vert}}f_{\text{inc},b}^{l+1/2}%
P_{b,\text{inc}} \label{dlh1}%
\end{align}

A familiar expression for the LH QL operator is obtained in cylindrical
geometry. From (\ref{divcyl}) and (\ref{dcyl1}-\ref{dcyl4}) we see that for
LHCD, the QL\ operator can be rewritten as
\begin{align}
Q^{\text{LH}}(f)  &  =\frac{\partial}{\partial p_{\Vert}}D_{\Vert\Vert
}^{\text{RF}}\frac{\partial f}{\partial p_{\Vert}}\nonumber\\
&  =\sum_{b}\frac{\partial}{\partial p_{\Vert}}\overline{D}_{b,0\text{,new}%
}^{\text{LH}}\frac{v_{Te}}{\left|  v_{\Vert}\right|  }H\left(  \frac
{c}{v_{\Vert}}-N_{b\Vert\min}\right)  H\left(  N_{b\Vert\max}-\frac
{c}{v_{\Vert}}\right)  \frac{\partial f}{\partial p_{\Vert}} \label{dlh}%
\end{align}

Although it is better to keep the exact formula (\ref{dlh}), the factor
$v_{Te}/\left|  v_{\Vert}\right|  $ has often been neglected in the
litterature, which is an acceptable approximation when $\overline
{D}_{b,0\text{,new}}^{\text{LH}}\gg\nu_{e}p_{Te}^{2}$ and $\Delta N_{b\Vert
}\ll N_{b\Vert\min}$. In this case, considering only one ray $b$, the LH
Operator reduces to
\begin{equation}
Q^{\text{LH}}(f)\simeq\left\{
\begin{array}
[c]{cc}%
\dfrac{\partial}{\partial p_{\Vert}}\overline{D}_{0\text{,new}}^{\text{LH}%
}\dfrac{\partial f}{\partial p_{\Vert}} & \text{for }v_{1}<\dfrac{v_{\Vert}%
}{v_{Te}}<v_{2}\\
0 & \text{otherwise}%
\end{array}
\right.
\end{equation}
with
\begin{align}
v_{1}  &  =\frac{1}{\beta_{Te}N_{\Vert\max}}\\
v_{2}  &  =\frac{1}{\beta_{Te}N_{\Vert\min}}\\
\overline{D}_{0\text{,new}}^{\text{LH}}  &  =\overline{\left[  \frac{v_{Te}%
}{\left|  v_{\Vert}\right|  }\right]  }\overline{D}_{b,0\text{,new}%
}^{\text{LH}}\\
&  =\frac{1}{rR_{p}}\frac{1}{\left|  \left(  v_{1}+v_{2}\right)  /2\right|
}\frac{1}{m_{e}\ln\Lambda}\frac{\omega}{\omega_{pe}^{4}}\frac{N_{\parallel,0}%
}{\Delta N_{\Vert}}f_{\text{inc}}^{l+1/2}P_{\text{inc}}%
\end{align}
where $\overline{\left[  \dfrac{v_{Te}}{\left|  v_{\Vert}\right|  }\right]  }$
is an averaged value of $\dfrac{v_{Te}}{\left|  v_{\Vert}\right|  }$ which can
be taken to be $\dfrac{1}{\beta_{Te}N_{\Vert,0}}$ and then%
\begin{equation}
\overline{D}_{0\text{,new}}^{\text{LH}}=\frac{1}{rR_{p}}\frac{1}{\beta_{Te}%
}\frac{1}{m_{e}\ln\Lambda}\frac{\omega}{\omega_{pe}^{4}}\frac{1}{\Delta
N_{\Vert}}f_{\text{inc}}^{l+1/2}P_{\text{inc}}%
\end{equation}

\section{Electron Cyclotron Current Drive\label{ECCD}}

The cold plasma description is usually a good approximation to determine the
ECW properties, as long as we stay away from the upper-hybrid resonance, where
mode-conversion to EBW occurs. However, even in the cold plasma model, the
polarizations are usually mixed for oblique propagation, and no simple
analytical formulation is possible. Still, it is possible to find limit cases
(small FLR effects, small $N_{\Vert}$) for which a simple analytical
derivation is possible.

\subsection{Polarization}

In the case of mostly perpendicular propagation, where $\left|  N_{\Vert
}\right|  \ll1$, the two modes are the quasi-X and quasi-O modes.

The polarization is mostly right-hand circular for X and parallel linear for
O. The only exception is for the X mode near the first harmonic $n=1$ where
$e_{b,-}^{\text{X}}$ $\simeq0$. For this reason, the first-harmonic is almost
transparent to the X-mode and this resonance is usually not considered.
Moreover, this harmonic can only be reached from the high field side, because
of the right-hand cut-off. From now, we consider only the X mode with $n\geq2$.

\subsection{Determination of $\Theta_{\mathbf{k}}^{b,\text{EC}}$}

\subsubsection{Small FLR\ limit}

Electron Cyclotron Current Drive (ECCD) results from momentum exchange from
the EC wave to the plasma through electron cyclotron damping at some harmonic
$n$. For a ray $b$, considering a given harmonic $n$, the coefficient
(\ref{thetakn2}) is
\begin{equation}
\Theta_{\mathbf{k}}^{b,n}=\frac{1}{\sqrt{2}}e_{b,+}e^{-i\alpha}J_{n-1}\left(
z_{b}\right)  +\frac{1}{\sqrt{2}}e_{b,-}e^{+i\alpha}J_{n+1}\left(
z_{b}\right)  +\frac{p_{\Vert}}{p_{\bot}}e_{b,\Vert}J_{n}\left(  z_{b}\right)
\end{equation}
where $z_{b}=k_{b\bot}v_{\bot}/\Omega$.

Since the resonance condition is $\omega_{b}-k_{b\Vert}v_{\Vert}-n\Omega=0$
and
\begin{equation}
\Omega=\frac{q_{e}B}{\gamma m_{e}}=-\frac{\omega_{ce}}{\gamma}<0
\end{equation}
the only harmonics to be considered for electrons are for
\begin{equation}
n\leq-1
\end{equation}

Applying the substitution
\begin{equation}
n^{\prime}=-n\geq1
\end{equation}
and renaming $n^{\prime}\rightarrow n\geq0$, we get
\begin{align}
\Theta_{\mathbf{k}}^{b,\text{EC}n} &  =\frac{1}{\sqrt{2}}e_{b,+}e^{-i\alpha
}J_{-n-1}\left(  z_{b}\right)  +\frac{1}{\sqrt{2}}e_{b,-}e^{+i\alpha}%
J_{-n+1}\left(  z_{b}\right)  +\frac{p_{\Vert}}{p_{\bot}}e_{b,\Vert}%
J_{-n}\left(  z_{b}\right) \nonumber\\
&  =\frac{1}{\sqrt{2}}e_{b,+}e^{-i\alpha}\left(  -1\right)  ^{n+1}%
J_{n+1}\left(  z_{b}\right)  +\frac{1}{\sqrt{2}}e_{b,-}e^{+i\alpha}\left(
-1\right)  ^{n-1}J_{n-1}\left(  z_{b}\right) \nonumber\\&\quad +\frac{p_{\Vert}}{p_{\bot}%
}e_{b,\Vert}\left(  -1\right)  ^{n}J_{n}\left(  z_{b}\right)
\label{thetaEC}
\end{align}

We have
\begin{equation}
z_{b}=\frac{k_{b\bot}v_{\bot}}{\Omega}=-N_{b\bot}\frac{p_{\bot}}{m_{e}c}%
\frac{\omega_{b}}{\omega_{ce}}%
\end{equation}
Using $\omega_{b}\sim n\omega_{ce}$, we get an estimation for $\left|
z_{b}\right|  $%
\begin{equation}
\left|  z_{b}\right|  \sim nN_{b\bot}\frac{p_{\bot}}{p_{Te}}\beta_{Te}%
\end{equation}
Typically, for ECCD in tokamaks, we have
\begin{equation}
N_{b\bot}\lesssim1,\;\beta_{Te}\lesssim0.1
\end{equation}
and most electrons of concern have
\begin{equation}
\frac{p_{\bot}}{p_{Te}}\leq3
\end{equation}
so that for low harmonics, the condition $\left|  z_{b}\right|  \ll1$ is
satisfied. In other words, the Larmor radius remains small compared to the
perpendicular wavelength. This condition is consistent with the conditions of
the cold plasma description of the EC waves. From now on, we assume $\left|
z_{b}\right|  \ll1$, which is the : limit of small FLR effects.

In this case, for $n\geq0$, the following approximate expression can be used
\begin{equation}
J_{n}\left(  z\right)  \simeq\frac{1}{n!}\left(  \frac{z}{2}\right)  ^{n}%
\end{equation}
and we have
\begin{align}
\Theta_{\mathbf{k}}^{b,\text{EC}n}\simeq&\frac{e_{b,+}}{\sqrt{2}}e^{-i\alpha
}\left(  -1\right)  ^{n+1}\frac{1}{\left(  n+1\right)  !}\left(  \frac{z_{b}%
}{2}\right)  ^{n+1}+\frac{e_{b,-}}{\sqrt{2}}e^{+i\alpha}\left(  -1\right)
^{n-1}\frac{1}{\left(  n-1\right)  !}\left(  \frac{z_{b}}{2}\right)
^{n-1}\nonumber\\&\ +\frac{p_{\Vert}}{p_{\bot}}e_{b,\Vert}\frac{1}{n!}\left(  \frac{z_{b}%
}{2}\right)  ^{n} \label{thetaklowz}%
\end{align}

\subsubsection{X-mode, $n\geq2$}

Because $e_{b,-}$ is the dominant polarization component and $\left|
z_{b}\right|  \ll1$, the second term in the sum in (\ref{thetaklowz}) is
largely dominant. We obtain%
\begin{equation}
\Theta_{\mathbf{k}}^{b,\text{EC-Xn}}\simeq e^{+i\alpha}\frac{e_{b,-}}{\sqrt
{2}}\frac{1}{\left(  n-1\right)  !}\left(  \frac{1}{2}N_{b\bot}\frac{p_{\bot}%
}{m_{e}c}\frac{\omega_{b}}{\omega_{ce}}\right)  ^{n-1} \label{thetakX}%
\end{equation}

\subsubsection{O-mode, $n\geq1$}

\bigskip For the O-mode, things are more complicated. However, if $\left|
N_{\Vert}\right|  \ll1$, $e_{b,\Vert}$ is much larger than $e_{b,-}$. The last
term in (\ref{thetaklowz}) is dominant if%
\begin{equation}
\left|  \frac{e_{b,-}}{e_{b,\Vert}}\right|  \ll\frac{1}{\sqrt{2}}\frac{\left|
p_{\Vert}\right|  }{m_{e}c}%
\end{equation}
which is satisfied for resonant electrons as long as the temperature is not
too low. In that case, we have%
\begin{equation}
\Theta_{\mathbf{k}}^{b,\text{EC-On}}\simeq\frac{p_{\Vert}}{p_{\bot}}%
e_{b,\Vert}\frac{1}{n!}\left(  \frac{1}{2}N_{b\bot}\frac{p_{\bot}}{m_{e}%
c}\frac{\omega_{b}}{\omega_{ce}}\right)  ^{n} \label{thetakO}%
\end{equation}

\subsection{Determination of $\mathbf{\Phi}_{b}^{\text{EC}}$ in the low
density limit.}

The vector $\mathbf{\Phi}_{b}$ describes the relation between the energy flux
and the electric field. In the cold plasma description, it is given by
(\ref{phibp2}).
\begin{equation}
\mathbf{\Phi}_{bP}=\text{Re}\left[  \mathbf{N}_{b}-\left(  \mathbf{N}_{b}%
\cdot\mathbf{e}_{b}\right)  \mathbf{e}_{b}^{\ast}\right]
\end{equation}

In general, both terms must be kept. However, where the density is low
\begin{equation}
\omega_{pe}\ll\omega_{b}%
\end{equation}
the ECW is mostly electromagnetic and mostly keeps its free-space
characteristics. In that case,%
\begin{align}
N_{b}  &  \simeq1\label{nb}\\
\frac{\mathbf{N}_{b}\cdot\mathbf{e}_{b}}{N_{b}}  &  \ll1
\end{align}
so that
\begin{equation}
\mathbf{\Phi}_{b}^{\text{EC}}\simeq\widehat{N}_{b} \label{phibec}%
\end{equation}

\subsection{EC Diffusion Coefficient}

\subsubsection{General expression with small FLR, $\left|  N_{\Vert}\right|
\ll1$ and $\omega_{pe}\ll\omega_{b}$ - or EM - approximations}

The normalized bounce-averaged diffusion coefficient for the Fokker-Planck
equation is given by (\ref{dbnb})%
\begin{align}
\overline{D}_{b,n}^{\text{EC}(0)}(p,\xi_{0}) &  =\frac{\gamma p_{Te}}{p\left|
\xi_{0}\right|  }\frac{1}{\lambda\widetilde{q}}\frac{r_{\theta_{b}}}{R_{p}%
}\frac{B^{\theta_{b}}}{B_{P}^{\theta_{b}}}\frac{\xi_{0}^{2}}{\xi_{\theta_{b}%
}^{2}}\Psi_{\theta_{b}}\overline{D}_{b,n,0}^{\text{EC,}\theta_{b}}%
\times\nonumber\\
&  H\left(  \theta_{b}-\theta_{\min}\right)  H\left(  \theta_{\max}-\theta
_{b}\right)  \left[  \frac{1}{2}\sum\limits_{\sigma}\right]  _{T}\delta\left(
N_{b\Vert}-N_{\Vert\text{res}}^{\theta_{b}}\right)  \left|  \Theta
_{\mathbf{k,}\theta_{b}}^{b,\left(  n\right)  }\right|  ^{2}\label{dec}%
\end{align}
with
\begin{align}
\overline{D}_{b,n,0}^{\text{EC,}\theta_{b}} &  =\frac{1}{r_{\theta_{b}%
}R_{\theta_{b}}}\frac{1}{m_{e}\ln\Lambda}\frac{1}{\omega_{b}\omega_{pe}^{2}%
}\frac{f_{\text{inc},b}^{l+1/2}}{\left|  \mathbf{\Phi}_{b}\right|
}P_{b,\text{inc}}\\
N_{\Vert\text{res}}^{\theta_{b}} &  =\frac{1}{\beta_{Te}}\frac{p_{Te}}%
{p\xi_{\theta_{b}}}\left(  \gamma-\frac{n\Psi_{\theta_{b}}\omega_{ce,0}%
}{\omega_{b}}\right)
\end{align}
where we used the substitution $n^{\prime}\rightarrow n$ as in (\ref{thetaEC}%
).

Within the small FLR, $\left|  N_{\Vert}\right|  \ll1$ and $\omega_{pe}%
\ll\omega_{b}$ - or electromagnetic - approximations, we have obtained the
following expressions for the ECW properties (\ref{thetakX}), ( \ref{thetakO}%
), (\ref{phibec})
\begin{align}
\Theta_{\mathbf{k}}^{b,\text{EC-Xn}} &  \simeq e^{+i\alpha}\frac{e_{b,-}%
}{\sqrt{2}}\frac{1}{\left(  n-1\right)  !}\left(  \frac{1}{2}N_{b\bot}%
\frac{p\sqrt{\Psi_{\theta_{b}}\left(  1-\xi_{0}^{2}\right)  }}{m_{e}c}%
\frac{\omega_{b}}{\Psi_{\theta_{b}}\omega_{ce,0}}\right)  ^{n-1}\\
\Theta_{\mathbf{k}}^{b,\text{EC-On}} &  \simeq\frac{\xi_{\theta_{b}}}%
{\sqrt{\Psi_{\theta_{b}}\left(  1-\xi_{0}^{2}\right)  }}e_{b,\Vert}\frac
{1}{n!}\left(  \frac{1}{2}N_{b\bot}\frac{p\sqrt{\Psi_{\theta_{b}}\left(
1-\xi_{0}^{2}\right)  }}{m_{e}c}\frac{\omega_{b}}{\Psi_{\theta_{b}}%
\omega_{ce,0}}\right)  ^{n}\\
\mathbf{\Phi}_{b}^{\text{EC}} &  \simeq\widehat{N}_{b}%
\end{align}
so that%
\[
\left|  \mathbf{\Phi}_{b}^{\text{EC}}\right|  \simeq1
\]

In addition, within these approximations, the polarizations take the following
limit expressions:%
\begin{gather}
\left|  e_{b,-}\right|  \simeq\frac{\sqrt{2}}{2}\text{ for the X mode }\left(
n\geq2\right)  \\
\left|  e_{b,\Vert}\right|  \simeq1\text{ for the O mode}%
\end{gather}

\subsubsection{Simplified expression for ECCD in the case of circular
concentric flux-surfaces}

We consider the limit of a tokamak with circular flux-surfaces. In that case,
we have the following identities
\begin{gather}
\left|  \widehat{\psi}\cdot\widehat{r}\right|  =1\\
\widetilde{q}=\frac{r}{R_{p}}\frac{B}{B_{P}}=\frac{r_{\theta_{b}}}{R_{p}}%
\frac{B^{\theta_{b}}}{B_{P}^{\theta_{b}}}%
\end{gather}
and the QL diffusion coefficients for ECCD\ (\ref{dec}) can therefore be
written as%
\begin{align}
\overline{D}_{b}^{\text{EC-Xn}(0)}(p,\xi_{0}) &  =\frac{\gamma p_{Te}%
}{p\left|  \xi_{0}\right|  }\frac{1}{\lambda}\frac{\xi_{0}^{2}}{\xi
_{\theta_{b}}^{2}}\Psi_{\theta_{b}}\overline{D}_{b,n,0}^{\text{EC,}\theta_{b}%
}\frac{1}{4\left[  \left(  n-1\right)  !\right]  ^{2}}\nonumber\\&\times\quad \left(  \frac{p}{p_{Te}%
}\right)  ^{2\left(  n-1\right)  }\left(
\frac{1-\xi_{0}^{2}}{\Psi _{\theta_{b}}}\right)  ^{\left(
n-1\right)  }\left(  \frac{1}{2}N_{b\bot
}\beta_{Te}\frac{\omega_{b}}{\omega_{ce,0}}\right)  ^{2\left(
n-1\right) }\nonumber\\&\times\quad   H\left(
\theta_{b}-\theta_{\min}\right)  H\left(  \theta_{\max}-\theta
_{b}\right)  \frac{1}{\sqrt{\pi}\Delta N_{\Vert}}\left[  \frac{1}{2}%
\sum\limits_{\sigma}\right]  _{T}\exp\left[  \frac{-\left(  N_{\Vert
\text{res}}-N_{b\Vert,0}\right)  ^{2}}{\Delta N_{\Vert}^{2}}\right]
\end{align}%
\begin{align}
\overline{D}_{b}^{\text{EC-On}(0)}(p,\xi_{0}) &  =\frac{\gamma p_{Te}%
}{p\left|  \xi_{0}\right|
}\frac{1}{\lambda}\frac{\xi_{0}^{2}}{\left( 1-\xi_{0}^{2}\right)
}\overline{D}_{b,n,0}^{\text{EC,}\theta_{b}}\frac {1}{\left[
n!\right]  ^{2}}\nonumber\\&\times\quad \left(
\frac{p}{p_{Te}}\right)  ^{2n}\left(
\frac{1-\xi_{0}^{2}}{\Psi_{\theta_{b}}}\right)  ^{n}\left(  \frac{1}%
{2}N_{b\bot}\beta_{Te}\frac{\omega_{b}}{\omega_{ce,0}}\right)  ^{2n}%
\nonumber\\&\times\quad  H\left( \theta_{b}-\theta_{\min}\right)
H\left(  \theta_{\max}-\theta
_{b}\right)  \frac{1}{\sqrt{\pi}\Delta N_{\Vert}}\left[  \frac{1}{2}%
\sum\limits_{\sigma}\right]  _{T}\exp\left[  \frac{-\left(  N_{\Vert
\text{res}}-N_{b\Vert,0}\right)  ^{2}}{\Delta N_{\Vert}^{2}}\right]
\end{align}
with
\begin{align}
\overline{D}_{b,n,0}^{\text{EC,}\theta_{b}} &  =\frac{1}{rR_{\theta_{b}}}%
\frac{1}{m_{e}\ln\Lambda}\frac{1}{\omega_{b}\omega_{pe}^{2}}f_{\text{inc}%
,b}^{l+1/2}P_{b,\text{inc}}\\
N_{\Vert\text{res}}^{\theta_{b}} &  =\frac{1}{\beta_{Te}}\frac{p_{Te}}%
{p\xi_{\theta_{b}}}\left(  \gamma-\frac{n\Psi_{\theta_{b}}\omega_{ce,0}%
}{\omega_{b}}\right)
\end{align}
and where we assume a gaussian power spectrum as in (\ref{gaussian}).

In order to compare with ECCD\ operators found in the litterature, we redefine
the EC constant factors such that%
\begin{align}
\overline{D}_{b}^{\text{EC-Xn}(0)}(p,\xi_{0}) &  =\frac{\gamma p_{Te}%
}{p\left|  \xi_{0}\right|  }\left(  \frac{p}{p_{Te}}\right)  ^{2\left(
n-1\right)  }\frac{\xi_{0}^{2}}{\xi_{\theta_{b}}^{2}}\overline{D}%
_{b,0\text{,new}}^{\text{EC-Xn,}\theta_{b}}\frac{\left(  1-\xi_{0}^{2}\right)
^{\left(  n-1\right)  }}{\Psi_{\theta_{b}}^{\left(  n-2\right)  }}\frac
{1}{\lambda}\nonumber\\
&  H\left(  \theta_{b}-\theta_{\min}\right)  H\left(  \theta_{\max}-\theta
_{b}\right)  \left[  \frac{1}{2}\sum\limits_{\sigma}\right]  _{T}\exp\left[
\frac{-\left(  N_{\Vert\text{res}}-N_{b\Vert,0}\right)  ^{2}}{\Delta N_{\Vert
}^{2}}\right]
\end{align}
with
\begin{align}
\overline{D}_{b,0\text{,new}}^{\text{EC-Xn,}\theta_{b}} &  =\frac{1}{\sqrt
{\pi}\Delta N_{\Vert}}\frac{1}{4\left[  \left(  n-1\right)  !\right]  ^{2}%
}\left(  \frac{1}{2}N_{b\bot}\beta_{Te}\frac{\omega_{b}}{\omega_{ce,0}%
}\right)  ^{2\left(  n-1\right)  }\overline{D}_{b,n,0}^{\text{EC,}\theta_{b}%
}\\
&  =\frac{1}{rR_{\theta_{b}}}\frac{1}{m_{e}\ln\Lambda}\frac{1}{\omega
_{b}\omega_{pe}^{2}}\frac{1}{\sqrt{\pi}\Delta N_{\Vert}}\frac{1}{4\left[
\left(  n-1\right)  !\right]  ^{2}}\left(  \frac{1}{2}N_{b\bot}\beta_{Te}%
\frac{\omega_{b}}{\omega_{ce,0}}\right)  ^{2\left(  n-1\right)  }%
f_{\text{inc},b}^{l+1/2}P_{b,\text{inc}}%
\end{align}
and%
\begin{align}
\overline{D}_{b}^{\text{EC-On}(0)}(p,\xi_{0}) &  =\frac{\gamma p_{Te}%
}{p\left|  \xi_{0}\right|  }\left(  \frac{p}{p_{Te}}\right)  ^{2n}\frac
{\xi_{0}^{2}\left(  1-\xi_{0}^{2}\right)  ^{n-1}}{\Psi_{\theta_{b}}^{n}%
}\overline{D}_{b,0\text{,new}}^{\text{EC-On,}\theta_{b}}\frac{1}{\lambda
}\nonumber\\
&  H\left(  \theta_{b}-\theta_{\min}\right)  H\left(  \theta_{\max}-\theta
_{b}\right)  \left[  \frac{1}{2}\sum\limits_{\sigma}\right]  _{T}\exp\left[
\frac{-\left(  N_{\Vert\text{res}}-N_{b\Vert,0}\right)  ^{2}}{\Delta N_{\Vert
}^{2}}\right]
\end{align}
with
\begin{align}
\overline{D}_{b,0\text{,new}}^{\text{EC-On,}\theta_{b}} &  =\frac{1}{\sqrt
{\pi}\Delta N_{\Vert}}\frac{1}{\left[  n!\right]  ^{2}}\left(  \frac{1}%
{2}N_{b\bot}\beta_{Te}\frac{\omega_{b}}{\omega_{ce,0}}\right)  ^{2n}%
\overline{D}_{b,n,0}^{\text{EC,}\theta_{b}}\\
&  =\frac{1}{rR_{\theta_{b}}}\frac{1}{m_{e}\ln\Lambda}\frac{1}{\omega
_{b}\omega_{pe}^{2}}\frac{1}{\sqrt{\pi}\Delta N_{\Vert}}\frac{1}{\left[
n!\right]  ^{2}}\left(  \frac{1}{2}N_{b\bot}\beta_{Te}\frac{\omega_{b}}%
{\omega_{ce,0}}\right)  ^{2n}f_{\text{inc},b}^{l+1/2}P_{b,\text{inc}}%
\end{align}

The two most common ECCD scenarios in experiments are the ones with the
largest diffusion coefficient: X2 and O1. For these case, we find%

\begin{align}
\overline{D}_{b}^{\text{EC-X2}(0)}(p,\xi_{0}) &  =\frac{\gamma p_{Te}%
}{p\left|  \xi_{0}\right|  }\left(  \frac{p}{p_{Te}}\right)  ^{2}\frac{\xi
_{0}^{2}\left(  1-\xi_{0}^{2}\right)  }{\xi_{\theta_{b}}^{2}}\overline
{D}_{b,0\text{,new}}^{\text{EC-X2,}\theta_{b}}\frac{1}{\lambda}\nonumber\\
&  H\left(  \theta_{b}-\theta_{\min}\right)  H\left(  \theta_{\max}-\theta
_{b}\right)  \left[  \frac{1}{2}\sum\limits_{\sigma}\right]  _{T}\exp\left[
\frac{-\left(  N_{\Vert\text{res}}-N_{b\Vert,0}\right)  ^{2}}{\Delta N_{\Vert
}^{2}}\right]
\end{align}
with
\begin{equation}
\overline{D}_{b,0\text{,new}}^{\text{EC-X2,}\theta_{b}}=\frac{1}%
{rR_{\theta_{b}}}\frac{1}{m_{e}\ln\Lambda}\frac{1}{\omega_{b}\omega_{pe}^{2}%
}\frac{1}{\sqrt{\pi}\Delta N_{\Vert}}\frac{1}{4}\left(  \frac{1}{2}N_{b\bot
}\beta_{Te}\frac{\omega_{b}}{\omega_{ce,0}}\right)  ^{2}f_{\text{inc}%
,b}^{l+1/2}P_{b,\text{inc}}%
\end{equation}
and%
\begin{align}
\overline{D}_{b}^{\text{EC-O1}(0)}(p,\xi_{0}) &  =\frac{\gamma p_{Te}%
}{p\left|  \xi_{0}\right|  }\left(  \frac{p}{p_{Te}}\right)  ^{2}\frac{\xi
_{0}^{2}}{\Psi_{\theta_{b}}}\overline{D}_{b,0\text{,new}}^{\text{EC-O1,}%
\theta_{b}}\frac{1}{\lambda}\nonumber\\
&  H\left(  \theta_{b}-\theta_{\min}\right)  H\left(  \theta_{\max}-\theta
_{b}\right)  \left[  \frac{1}{2}\sum\limits_{\sigma}\right]  _{T}\exp\left[
\frac{-\left(  N_{\Vert\text{res}}-N_{b\Vert,0}\right)  ^{2}}{\Delta N_{\Vert
}^{2}}\right]
\end{align}
with
\begin{equation}
\overline{D}_{b,0\text{,new}}^{\text{EC-O1,}\theta_{b}}=\frac{1}%
{rR_{\theta_{b}}}\frac{1}{m_{e}\ln\Lambda}\frac{1}{\omega_{b}\omega_{pe}^{2}%
}\frac{1}{\sqrt{\pi}\Delta N_{\Vert}}\left(  \frac{1}{2}N_{b\bot}\beta
_{Te}\frac{\omega_{b}}{\omega_{ce,0}}\right)  ^{2}f_{\text{inc},b}%
^{l+1/2}P_{b,\text{inc}}%
\end{equation}
